% =====================================================================
%  Implementacion (7 diapositivas: motor y esquema hibrido, Guia 2.2)
% =====================================================================

% SLIDE impl_arquitectura | V:A | T:25
\begin{frame}{Un binario, tres esquemas}
  \centering
  \begin{tikzpicture}[scale=0.9,
      box/.style={draw,rounded corners=2pt,font=\scriptsize,align=center,minimum height=0.7cm,
                  inner sep=3pt},
      sch/.style={box,minimum width=3.4cm,thick},
      mod/.style={box,minimum width=3.0cm,fill=gray!12},
      arr/.style={-{Latex[length=2mm]},gray!70!black}]
    \node[box,minimum width=3.6cm] (main) at (0,1.9) {\texttt{main} \texttt{-{}-scheme}};
    \node[sch,draw=tsdmdc,text=tsdmdc] (ts) at (-4.2,0.4) {\texttt{runEvacuation}\\TSDMD};
    \node[sch,draw=edmdc,text=edmdc] (ed) at (0,0.4) {\texttt{runEdmd}\\EDMD};
    \node[sch,draw=hybc,text=hybc] (hy) at (4.2,0.4) {\texttt{runHybrid}\\híbrido};
    \node[mod] (fo) at (-5.3,-1.3) {\texttt{forces}\\SFM y Cell Index Method};
    \node[mod] (ek) at (-1.7,-1.3) {\texttt{edmd\_kernel}\\tiempos de contacto};
    \node[mod] (re) at (1.7,-1.3) {\texttt{relaxation.h}\\recuperación de $v_d$};
    \node[mod] (ge) at (5.3,-1.3) {\texttt{generator}\\condición inicial};
    \node[box,minimum width=6cm] (out) at (0,-2.6) {salida de texto: un frame cada $dt_2$};
    \draw[arr] (main) -- (ts);
    \draw[arr] (main) -- (ed);
    \draw[arr] (main) -- (hy);
    \draw[arr] (ts) -- (fo);
    \draw[arr] (ed) -- (ek);
    \draw[arr] (ed) -- (re);
    \begin{pgfonlayer}{background}\draw[arr,gray!50] (hy) -- (fo);\end{pgfonlayer}
    \draw[arr] (hy) -- (ek);
    \draw[arr] (hy) -- (re);
    \draw[arr] (fo.south) -- (fo.south |- out.north);
    \draw[arr] (ek.south) -- (ek.south |- out.north);
    \draw[arr] (ge.south) -- (ge.south |- out.north);
  \end{tikzpicture}
  \vspace{2pt}

  {\small Un hilo, C++20, mismo \texttt{-O2} en los tres esquemas.}
\end{frame}

% SLIDE impl_tsdmd | V:A | T:30
\begin{frame}{TSDMD: Velocity Verlet con velocidad provisoria}
  \small
  \begin{align*}
    \vect x(t+dt) &= \vect x(t) + \vect v(t)\,dt + \tfrac12\,\vect a(t)\,dt^2\\
    \vect a(t+dt) &= \vect a\bigl(\vect x(t+dt),\,\vect v(t)+\vect a(t)\,dt\bigr)\\
    \vect v(t+dt) &= \vect v(t) + \tfrac12\bigl[\vect a(t)+\vect a(t+dt)\bigr]\,dt
  \end{align*}
  \begin{itemize}
    \item $dt = 10^{-4}\un{s}$, por debajo de la cota $0.1\sqrt{m/k}=2.58\times10^{-3}\un{s}$.
    \item Vecinos con Cell Index Method en cada paso.
    \item Cruce de la puerta (segundo objetivo a $1\un{m}$) y retiro del agente.
  \end{itemize}
\end{frame}

% SLIDE impl_edmd | V:A | T:35
\begin{frame}{EDMD: lista de eventos}
  \small
  Trayectoria libre exacta entre eventos:
  \[
    \vect x(t) = \vect x_0 + v_d\,\hat{\vect e}\,t
      + (\vect v_0 - v_d\,\hat{\vect e})\,\tau\,(1-e^{-t/\tau})
  \]
  \begin{itemize}
    \item Eventos: contacto agente-agente, con pared y con borde de puerta; cruce de
      puerta; retiro a $1\un{m}$; cambio de celda; \texttt{Steer} cada $0.1\un{s}$.
    \item Tiempo de contacto: raíz de $|\Delta\vect x(t)| = r_i + r_j$, bisección
      certificada a $10^{-13}\un{s}$.
    \item Cola binaria ordenada por (tiempo, secuencia).
    \item \textbf{Contador de generación} por agente: un evento con generación vieja
      se descarta como obsoleto.
  \end{itemize}
\end{frame}

% SLIDE impl_colision | V:A | T:35
\begin{frame}{Colisión con restitución y recuperación de $v_d$}
  \small
  \begin{columns}[T]
    \begin{column}{0.54\textwidth}
      Restitución $\varepsilon = 0.5$ solo en la normal $\hat{\vect n}$:
      \begin{align*}
        \vect v_i' &= \vect v_i - \tfrac{1+\varepsilon}{2}\,[(\vect v_i-\vect v_j)\cdot\hat{\vect n}]\,\hat{\vect n}\\
        \vect v_j' &= \vect v_j + \tfrac{1+\varepsilon}{2}\,[(\vect v_i-\vect v_j)\cdot\hat{\vect n}]\,\hat{\vect n}
      \end{align*}
      \begin{itemize}
        \item La parte tangencial no cambia; se conserva el momento del par.
        \item Pared: plano fijo. Borde de puerta: punto fijo.
        \item Se separa el solape residual; piso \texttt{event-gap} de $5\un{ms}$ por contacto
          contra el colapso inelástico.
      \end{itemize}
    \end{column}
    \begin{column}{0.42\textwidth}
      Recuperación de la velocidad deseada tras la colisión:
      \[
        \vect v(t) = v_d\,\hat{\vect e} + (\vect v_0 - v_d\,\hat{\vect e})\,e^{-t/\tau}
      \]
      La misma ley en los tres esquemas: la da la dinámica en TSDMD y \texttt{relaxation.h}
      en EDMD y en el híbrido.
    \end{column}
  \end{columns}
\end{frame}

% SLIDE hyb_criterio | V:B | T:50
\begin{frame}{Híbrido: criterio de pasaje}
  \small
  \begin{columns}[T]
    \begin{column}{0.50\textwidth}
      \[
        \tau_c(i) = \min\{t_{ij},\,t_{iw},\,t_{ic}\}
      \]
      Tiempo hasta el próximo contacto (vecinos, paredes, bordes de puerta), con el kernel
      de EDMD; el vecino TSDMD se extrapola a velocidad constante.
      \begin{itemize}
        \item EDMD $\to$ TSDMD: $\tau_c < T$.
        \item TSDMD $\to$ EDMD: $\tau_c > \alpha T$, sin contacto ni solape vigente y sin colapso.
        \item Banda $[T,\alpha T]$: histéresis, no hay cambio.
        \item $T = 0.5\un{s}$, $\alpha = 2$, $T \ge 2\,dt$.
      \end{itemize}
    \end{column}
    \begin{column}{0.46\textwidth}
      \centering
      \begin{tikzpicture}[
          st/.style={draw,thick,rounded corners=3pt,minimum width=1.8cm,minimum height=0.7cm,
                     font=\scriptsize\bfseries},
          arr/.style={-{Latex[length=2mm]},thick}]
        \node[st,draw=edmdc,text=edmdc] (e) at (-1.5,1.5) {EDMD};
        \node[st,draw=tsdmdc,text=tsdmdc] (t) at (1.5,1.5) {TSDMD};
        \draw[arr] (e.north east) to[bend left=35] node[above,font=\scriptsize] {$\tau_c<T$} (t.north west);
        \draw[arr] (t.south west) to[bend left=35] node[below,font=\scriptsize] {$\tau_c>\alpha T$} (e.south east);
        \draw[-{Latex[length=2mm]}] (-2.2,-0.5) -- (2.4,-0.5) node[right,font=\scriptsize] {$\tau_c$};
        \foreach \x/\l in {-1.2/0, 0/T, 1.2/\alpha T} {
          \draw (\x,-0.43) -- (\x,-0.57);
          \node[font=\scriptsize,below] at (\x,-0.57) {$\l$};
        }
        \fill[tsdmdc!25] (-1.2,-0.15) rectangle (0,0.15);
        \fill[gray!25] (0,-0.15) rectangle (1.2,0.15);
        \fill[edmdc!25] (1.2,-0.15) rectangle (2.2,0.15);
        \node[font=\tiny] at (-0.6,0.0) {TSDMD};
        \node[font=\tiny] at (0.6,0.0) {histéresis};
        \node[font=\tiny] at (1.7,0.0) {EDMD};
      \end{tikzpicture}

      \vspace{4pt}
      {\footnotesize $T=0$: EDMD puro. $T\ge t_{max}$: TSDMD puro. Trayectorias idénticas
      byte a byte.}
    \end{column}
  \end{columns}
\end{frame}

% SLIDE hyb_sync | V:B | T:55
\begin{frame}{Híbrido: un solo reloj}
  \small
  \begin{columns}[T]
    \begin{column}{0.54\textwidth}
      Un solo reloj. Grilla $t_k = k\,dt$ mientras haya algún agente TSDMD; sin ellos el reloj
      salta de evento en evento. Los eventos EDMD anteriores a $t_k$ se procesan entre límites.
      Cinco pasos en cada límite $t_k$:
      \begin{enumerate}\setlength{\itemsep}{0pt}
        \item un paso de Verlet del subconjunto TSDMD;
        \item el motor EDMD avanza hasta $t_k$;
        \item cruce de puerta y retiro;
        \item pasajes EDMD $\to$ TSDMD (eventos \texttt{Promote}, escaneo de pares mixtos);
        \item pasajes TSDMD $\to$ EDMD de a uno, en orden de id, reconstruyendo las fuerzas.
      \end{enumerate}
      Los pasajes ocurren solo en límites.
    \end{column}
    \begin{column}{0.42\textwidth}
      \centering
      \begin{tikzpicture}[ev/.style={circle,fill=edmdc,inner sep=1.6pt}]
        \draw[-{Latex[length=2mm]},thick] (-1.0,0) -- (3.9,0) node[right,font=\scriptsize] {$t$};
        \foreach \x/\l in {0/t_{k-1}, 1.5/t_k, 3.0/t_{k+1}} {
          \draw[thick] (\x,-0.2) -- (\x,0.2);
          \node[font=\scriptsize,below=3pt] at (\x,-0.2) {$\l$};
        }
        \node[ev] at (0.35,0) {};
        \node[ev] at (0.9,0) {};
        \node[ev] at (2.0,0) {};
        \node[ev] at (2.4,0) {};
        \node[ev] at (2.75,0) {};
        \node[font=\scriptsize,edmdc,above=3pt] at (0.6,0) {eventos EDMD};
        \draw[decorate,decoration={brace,amplitude=4pt}] (1.5,0.9) -- (0,0.9)
          node[midway,above=4pt,font=\scriptsize] {entre límites};
        \node[draw,rounded corners=2pt,font=\scriptsize,align=center,fill=gray!12]
          at (1.5,-1.9) {en $t_k$: pasos (1) a (5)};
        \draw[-{Latex[length=2mm]},gray!70!black] (1.5,-1.0) -- (1.5,-1.5);
      \end{tikzpicture}
    \end{column}
  \end{columns}
\end{frame}

% SLIDE hyb_contactos | V:B | T:45
\begin{frame}{Híbrido: contactos sin pérdidas}
  \small
  \begin{itemize}
    \item Residente EDMD hacia un agente TSDMD: evento \texttt{Promote} en $t_c - T$; el
      contacto lo resuelve el resorte SFM (0 colisiones duras en el híbrido).
    \item Residente EDMD que se acerca a un agente TSDMD, sin eventos propios: escaneo de
      pares mixtos en las celdas $3\times3$, con el vecino TSDMD a velocidad constante.
    \item El pasaje conserva $\vect x$ y $\vect v$, sube la generación del agente y reinicia
      la relajación con la velocidad actual.
    \item Auditoría por fuerza bruta en cada límite de una corrida completa: 0 solapes de
      residente EDMD, 0 contactos perdidos y 0 violaciones del oráculo, pares EDMD-TSDMD
      incluidos. Tres controles negativos se detectan.
  \end{itemize}
  \vfill
  {\footnotesize Alcance de la auditoría: $N=60$ y $N=200$ (semilla 1) y dos escenarios
  armados.}
\end{frame}
